\documentclass[11pt]{article}
\usepackage[margin=1in]{geometry}
\usepackage{amsmath,amssymb,amsthm}
\usepackage{algorithm}
\usepackage{algpseudocode}
\usepackage{booktabs}
\usepackage{hyperref}

\newtheorem{definition}{Definition}
\newtheorem{theorem}{Theorem}
\newtheorem{lemma}[theorem]{Lemma}
\newtheorem{property}{Property}

\title{\textbf{Mathematical Specification, Pathwise Derivatives, and Continuous Relaxations of Gaussian Reparameterization and Gumbel-Softmax (ALGO-NN-30)}}
\author{Team Scaibu \\ \small Email: \texttt{jaydeep.9664920749@gmail.com} \\ \small Architecture and Core Engine Team}
\date{\today}

\begin{document}
\maketitle

\begin{abstract}
Differentiating expectations of stochastic nodes $\nabla_\theta \mathbb{E}_{q_\theta(z)}[f(z)]$ using standard score-function estimators exhibits high variance that impairs optimization. Pathwise reparameterization reformulates continuous variables via independent base noise $z = \mu + \sigma \odot \epsilon$, while the Gumbel-Softmax (Concrete) distribution provides a continuous, differentiable surrogate for discrete categorical choices. We formalize both sampling mechanics, prove the Low Variance Gradient Theorem, provide full pseudocode matching the reference implementation, and derive complexity.
\end{abstract}

\section{Notation and Mathematical Preliminaries}
\label{sec:notation}

Let $\boldsymbol{\mu} \in \mathbb{R}^D$ denote the latent mean, $\log\boldsymbol{\sigma}^2 \in \mathbb{R}^D$ the latent log-variance, and $\boldsymbol{\epsilon} \sim \mathcal{N}(\mathbf{0}, \mathbf{I}_D)$ standard normal noise. For categorical distributions, let $\mathbf{z} \in \mathbb{R}^K$ denote unnormalized logits and $g_k \sim \operatorname{Gumbel}(0, 1)$ standard Gumbel noise generated from uniform variables $u_k \sim \mathcal{U}(0, 1)$ via $g_k = -\log(-\log u_k)$.

\subsection{Pathwise Continuous and Discrete Relaxations}
\paragraph{Intuitive Meaning:} By expressing a random sample as a deterministic formula of distribution parameters plus independent background noise, gradients can flow directly through the sample into the parameters using standard backpropagation.

\paragraph{Formal Mathematical Definition:}
\begin{itemize}
    \item \textbf{Gaussian Reparameterization:}
    \begin{equation}
    \label{eq:gaussian_reparam}
    \mathbf{z} = g(\boldsymbol{\mu}, \boldsymbol{\sigma}, \boldsymbol{\epsilon}) = \boldsymbol{\mu} + \boldsymbol{\sigma} \odot \boldsymbol{\epsilon} = \boldsymbol{\mu} + \exp\left(\frac{1}{2}\log\boldsymbol{\sigma}^2\right) \odot \boldsymbol{\epsilon}.
    \end{equation}
    \item \textbf{Gumbel-Softmax Continuous Relaxation:} For temperature $\tau > 0$:
    \begin{equation}
    \label{eq:gumbel_softmax}
    y_k = \frac{\exp((z_k + g_k)/\tau)}{\sum_{j=1}^K \exp((z_j + g_j)/\tau)}.
    \end{equation}
\end{itemize}

\paragraph{Worked Numerical Example:}
Let $\mu = 1.0, \log\sigma^2 = 0.0 \implies \sigma = 1.0, \epsilon = 0.5$:
$z = 1.0 + 1.0(0.5) = 1.5$.
Gradients: $\frac{\partial z}{\partial \mu} = 1.0, \frac{\partial z}{\partial \log\sigma^2} = 0.5 \sigma \epsilon = 0.5(1.0)(0.5) = 0.25$.

\subsection{Summary of Symbols}
\begin{itemize}
    \item $D \in \mathbb{N}$: Dimension of continuous latent space.
    \item $K \in \mathbb{N}$: Number of discrete categorical classes.
    \item $\boldsymbol{\mu} \in \mathbb{R}^D$: Latent Gaussian mean.
    \item $\log\boldsymbol{\sigma}^2 \in \mathbb{R}^D$: Latent Gaussian log-variance.
    \item $\boldsymbol{\epsilon} \in \mathbb{R}^D$: Standard normal noise vector.
    \item $\mathbf{z} \in \mathbb{R}^K$: Unnormalized categorical logit vector.
    \item $\mathbf{g} \in \mathbb{R}^K$: Gumbel noise vector.
    \item $\tau \in \mathbb{R}_{>0}$: Softmax temperature parameter.
    \item $\mathbf{y} \in \Delta^{K-1}$: Differentiable soft one-hot probability vector.
\end{itemize}

\section{Problem Definition and Core Concepts}
\label{sec:problem_definition}

\subsection{Temperature Annealing Dynamics}
As temperature $\tau \to 0$, the Gumbel-Softmax sample $\mathbf{y}$ converges almost surely to an exact discrete one-hot vector $\operatorname{argmax}_k(z_k + g_k)$. As $\tau \to \infty$, $\mathbf{y}$ approaches the uniform distribution $\left[\frac{1}{K}, \dots, \frac{1}{K}\right]$.

\section{Algorithmic Specification}
\label{sec:algorithm}

Algorithm~\ref{alg:reparam_gumbel} specifies the sampling procedure matching \texttt{impl.py}.

\begin{algorithm}[H]
\caption{Continuous and Discrete Reparameterization Sampling}
\label{alg:reparam_gumbel}
\begin{algorithmic}[1]
\Require Mode $\mathrm{mode} \in \{\text{gaussian}, \text{gumbel\_softmax}\}$, parameters $(\boldsymbol{\mu}, \log\boldsymbol{\sigma}^2, \boldsymbol{\epsilon})$ or $(\mathbf{z}, \mathbf{g}, \tau, \mathrm{hard})$.
\Ensure Sample vector $\mathbf{s}$.
\If{$\mathrm{mode} = \text{gaussian}$}
    \State Initialize $\mathbf{s} \gets []$
    \For{$i = 1$ to $D$}
        \State $\sigma_i \gets \exp(0.5 \cdot \log\sigma^2_i)$
        \State $s_i \gets \mu_i + \sigma_i \cdot \epsilon_i$
        \State Append $s_i$ to $\mathbf{s}$
    \EndFor
    \State \Return $\mathbf{s}$
\Else
    \State $\mathbf{p} \gets \left[ \frac{z_k + g_k}{\tau} \right]_{k=1}^K$
    \State $m \gets \max_k p_k$
    \State $S \gets \sum_{k=1}^K \exp(p_k - m)$
    \State $\mathbf{y} \gets \left[ \frac{\exp(p_k - m)}{S} \right]_{k=1}^K$
    \If{$\mathrm{hard} = \text{True}$}
        \State $k^* \gets \operatorname{argmax}_k y_k$
        \State $\mathbf{s} \gets [\mathbb{I}_{\{k = k^*\}}]_{k=1}^K$
    \Else
        \State $\mathbf{s} \gets \mathbf{y}$
    \EndIf
    \State \Return $\mathbf{s}$
\EndIf
\end{algorithmic}
\end{algorithm}

\section{Theoretical Guarantees and Variance Reduction}
\label{sec:guarantees}

\begin{theorem}[Pathwise Gradient Variance Reduction]
\label{thm:pathwise_variance}
\textbf{Assumptions:} Let $f(z)$ be $L$-Lipschitz continuous, and $z \sim \mathcal{N}(\mu, \sigma^2)$.
\textbf{Guarantee:} The pathwise gradient estimator $\hat{g}_{\mathrm{path}} = \nabla_z f(g(\mu, \sigma, \epsilon)) \nabla_\mu g = \nabla_z f(\mu + \sigma \epsilon)$ has finite variance $\operatorname{Var}(\hat{g}_{\mathrm{path}}) \le L^2$, whereas score-function (REINFORCE) variance scales with the magnitude of $f(z)^2 / \sigma^2$.
\textbf{Proof:}
Because $f$ is $L$-Lipschitz, $|\nabla_z f(z)| \le L$ almost everywhere.
$\operatorname{Var}(\hat{g}_{\mathrm{path}}) = \mathbb{E}[(\nabla_z f)^2] - (\mathbb{E}[\nabla_z f])^2 \le \mathbb{E}[L^2] - 0 = L^2$.
By contrast, score function $\hat{g}_{\mathrm{score}} = f(z) \frac{z - \mu}{\sigma^2}$ has variance $\mathbb{E}\left[ f(z)^2 \frac{(z-\mu)^2}{\sigma^4} \right] = \mathcal{O}\left(\frac{\|f\|_\infty^2}{\sigma^2}\right)$, which blows up as $\sigma \to 0$.
\textbf{Limitations:} Pathwise gradients require the mapping $f(z)$ to be differentiable with respect to sample $z$.
\end{theorem}

\section{Complexity Analysis}
\label{sec:complexity}

\begin{itemize}
    \item \textbf{Time Complexity:} $\mathcal{O}(D)$ operations for Gaussian sampling; $\mathcal{O}(K)$ operations for Gumbel-Softmax normalization.
    \item \textbf{Space Complexity:} $\mathcal{O}(D)$ or $\mathcal{O}(K)$ for sample storage.
\end{itemize}

\section{Worked Numerical Example}
\label{sec:worked_example}

Let $\mathbf{z} = [1.0, 0.0], \mathbf{g} = [0.5, -0.2], \tau = 1.0$:
\begin{enumerate}
    \item Perturbed logits: $\mathbf{p} = [1.0 + 0.5, 0.0 - 0.2] = [1.5, -0.2]$.
    \item Exponentials: $e^{1.5} = 4.481689, e^{-0.2} = 0.818731$.
    \item Sum: $4.481689 + 0.818731 = 5.300420$.
    \item Normalized sample: $\mathbf{y} = \left[\frac{4.481689}{5.300420}, \frac{0.818731}{5.300420}\right] = [0.845535, 0.154465]$.
\end{enumerate}

\section{Numerical Considerations and Edge Cases}
\label{sec:numerical}

\begin{itemize}
    \item \textbf{Straight-Through Gumbel Hard Mode:} In hard mode, the forward pass outputs discrete one-hot vector $[1.0, 0.0]$ while passing soft continuous probabilities $\mathbf{y}$ into the backward computational graph via STE.
\end{itemize}

\begin{thebibliography}{9}
\bibitem{kingma2013} D.~P.~Kingma and M.~Welling, ``Auto-Encoding Variational Bayes,'' \emph{ICLR}, 2014.
\bibitem{jang2016} E.~Jang, S.~Gu, and B.~Poole, ``Categorical Reparameterization with Gumbel-Softmax,'' \emph{ICLR}, 2017.
\end{thebibliography}

\end{document}
